All arithmetic claims are over the fixed-point integer encoding implemented by
the Python relation and Rust/SP1 backend. They are not claims about exact
real-valued PyTorch execution.

\paragraph{Merkle membership.}
The membership relation recomputes a canonical transition leaf hash, follows
the supplied sibling path and direction bits, and checks equality with the
public dataset root. Dataset-size scaling cases for 1k, 10k, and 100k leaves
are proof-verified.

\paragraph{Bellman target and TD MVP.}
The TD relation checks reward, done flag, action, discount factor, selected Q
value, target Q value, Bellman target, TD error, and SmoothL1 loss. In the TD
MVP vector the Q values are witnesses; in model-grounded relations they are
derived from committed tiny MLP weights.

\paragraph{Forward-TD MLP and SmoothL1.}
The Forward-TD MLP relation binds quantized online and target weights, runs a
Linear-ReLU-Linear forward pass, checks argmax selection, derives the selected
target value, and recomputes SmoothL1 loss.

\paragraph{Backpropagation and one-step SGD.}
The one-step training-update and one-step SGD tiny relations recompute the
SmoothL1 derivative branch, fixed-point gradient/backprop path, SGD deltas,
post-update weights, and next checkpoint hash. The supported update is simple
SGD for canonical tiny vectors, not Adam and not batch sizes 4/8/16.

\paragraph{Checkpoint and target-network relations.}
Checkpoint hashes bind model states across one-step updates, short traces, and
training fragments. Fragment relations additionally check deterministic
sampling, global steps, target checkpoint links, and hard target-network sync.

\paragraph{Proof-manifest chunk chains.}
The aggregation relation binds child proof-manifest hashes, child public-input
hashes, chunk order, step ranges, dataset/config consistency, aggregate roots,
and online/target checkpoint links. It does not recursively verify child proof
bytes inside SP1.
